% Experiment report. Fixed sections: see experiments/TEMPLATE/report.tex and experiments/README.md.
% Build: npm run experiment:report -- 13-rl-full-state-value
\documentclass{../polyexp}
\begin{document}
% generated by scripts/experiment-report.ts from experiment.json and models/<run>/ - do not edit
\providecommand{\ExpHeader}{}\renewcommand{\ExpHeader}{%
{\sffamily\footnotesize\color{expgrey}EXPERIMENT 13 \textperiodcentered\ METHOD \hfill 2026-11-25}\par
{\LARGE\bfseries\sffamily Full integer state: value function learned from recorded games}\hfill\verdictbox{expred}{NOT KEPT}\par\smallskip
{\itshape Can a value function over the exact integer state (turn, cities, income, units, sum of city levels, techs), learned from recorded Perfection games, predict the learner's remaining score better than the turn alone, and does a learned transition model (backward induction) predict better than Monte Carlo averages?}\par\medskip
\begin{tabular}{@{}p{22mm}p{\dimexpr\linewidth-26mm}@{}}
\fact{Change}{New tool engine/rlvalue (rl-value): records the six-integer state of seat 0 at the start of each of its turns, learns a tabular transition and reward model, computes V by backward induction from the last turn, and compares it with a Monte Carlo table and a turn-only baseline on held-out games.}
\fact{Base}{\model{greedy-full-probe}}
\fact{Method}{Training-set sizes 1k to 300k games (seeds 1..), 10,000 held-out games (seeds 500,000,000..), RMSE of the predicted return-to-go; metric: RMSE of predicted remaining score}
\fact{Setups}{Perfection, 16$\times$16 continents, 3 players, Crazy, seat 0 Xin-Xi (human seat), greedy-full-probe in every seat}
\fact{Record}{issue \issue{166}; status: done}
\end{tabular}\par\medskip
\begin{center}\begin{tikzpicture}[node distance=5mm]
    \node[flowbox] (n0) {{\sffamily\bfseries\scriptsize\color{expgrey}BASE}\\[1pt]\model{full-probe}};
    \node[flowbox,right=5mm of n0] (n1) {{\sffamily\bfseries\scriptsize\color{expgrey}CHANGE}\\[1pt]method};
    \node[flowbox,right=5mm of n1] (n2) {{\sffamily\bfseries\scriptsize\color{expgrey}MEASURED}\\[1pt]Held-out RMSE, 10,000 games\\{\scriptsize 1 setup, metric: RMSE of predicted remaining score}};
    \node[flowbox,right=5mm of n2] (n3) {{\sffamily\bfseries\scriptsize\color{expgrey}RESULT}\\[1pt]Held-out RMSE of remaining score, 300k training games: turn baseline vs model-based V\\{\scriptsize 2414.3 $\rightarrow$ 2210.4}\\\textbf{$-$203.9}};
    \node[flowbox,right=5mm of n3] (n4) {{\sffamily\bfseries\scriptsize\color{expgrey}VERDICT}\\[1pt]\textbf{\color{expred}NOT KEPT}};
    \draw[flowarrow] (n0) -- (n1); \draw[flowarrow] (n1) -- (n2); \draw[flowarrow] (n2) -- (n3); \draw[flowarrow] (n3) -- (n4);
\end{tikzpicture}\end{center}
\begin{center}\sffamily\small\begin{tabular}{@{}lrrr@{}}\toprule
Measure & Base & New & Change\\\midrule
Held-out RMSE of remaining score, 300k training games: turn baseline vs model-based V & 2414.3 & 2210.4 & $-$203.9\\
Same, turn baseline vs Monte Carlo table & 2414.3 & 2180.6 & $-$233.7\\
Turn 20 only: turn baseline vs model-based V & 2124.3 & 1595.0 & $-$529.3\\
\bottomrule\end{tabular}\end{center}
}
\providecommand{\ExpResults}{}\renewcommand{\ExpResults}{%
{\small\color{expgrey}No runs recorded.}
}

\ExpHeader

\section{Question and hypothesis}
% WHEN: before the first run.
\textbf{The idea (David, 25 Nov 2026, issue \#166).} Model the game state directly with all features at once,
as exact integers and without bins: number of cities, income (stars per turn), number of units, sum of city
levels, number of techs, and the turn (Perfection ends after 30 turns, so time is part of the state). Learn the
value of these states from play; later, use value iteration on them to play the game. This runs in parallel
with Jakob's track, which starts from cities alone and adds one feature at a time.

\textbf{This experiment} is the first step: the value function of the current bot's own play (policy
evaluation, no decisions yet). The value of a state is the expected remaining score,
$V(s_t) = \mathbb{E}[\,\text{final score} - \text{score}_t \mid s_t\,]$, which is the sum of the rewards
(score changes) still to come.

\textbf{Prediction} (written before the main run; a 200-game pilot was run only to check the tool and time it):
\begin{enumerate}
  \item \emph{The state space is sparse.} Distinct states grow almost in proportion to the games played. With
        300{,}000 training games, more than half of the distinct states are seen only once, and the held-out
        coverage (share of held-out visits whose exact state was seen in training) is about 85\%: close to
        100\% in the first turns, below 60\% in the last ten.
  \item \emph{The learned model beats Monte Carlo.} On held-out visits to states seen in training, the
        model-based $V$ (backward induction) has a lower RMSE than the Monte Carlo table at every training size,
        by at least 5\% at 300{,}000 games, because it pools the successors of rarely seen states.
  \item \emph{The full state beats the turn alone, given enough data.} At small sizes (the pilot: 200 games) the
        tables are worse than the turn-only baseline; from about 30{,}000 games on, model-based $V$ beats it,
        and at 300{,}000 games its RMSE over all held-out visits is at least 25\% below the baseline's.
\end{enumerate}
\textbf{Against the idea:} if model-based $V$ is not at least 10\% below the turn baseline at 300{,}000 games,
the exact-integer state is too sparse to learn as a table at this data size.

\section{Mechanism}
\textbf{State.} At the start of each turn of seat 0 (the learner), \texttt{rl-value} (\texttt{engine/rlvalue})
reads $s = (\text{turn}, \text{cities}, \text{income}, \text{units}, \text{levels}, \text{techs})$:
cities and units owned by seat 0, \texttt{state::income} (stars at the start of its next turn), the sum of
\texttt{level} over its cities, and the number of techs owned. It also records \texttt{score\_of}, the raw
Perfection score. A terminal row holds the final score.

\textbf{Reward.} $r_t = \text{score}_{t+1} - \text{score}_t$; the last step runs to the final score.

\textbf{Three estimators, all fitted on the same training games.}
\begin{itemize}
  \item \emph{Model-based (certainty equivalence).} For each visited state the list of observed transitions
        $(r, s')$ is the maximum-likelihood model $\hat P(s'\mid s)$, $\hat r(s,s')$. Backward induction from the
        last turn: $V(s) = \frac{1}{n(s)} \sum_{(r,s')} \bigl(r + V(s')\bigr)$, $V(\text{terminal}) = 0$. A
        successor is always one turn later, so this is exact finite-horizon dynamic programming, no
        discounting. It is also the fixed point that batch TD(0) converges to.
  \item \emph{Monte Carlo.} $V(s)$ = mean return-to-go of the training visits to $s$.
  \item \emph{Turn baseline.} $V(t)$ = mean return-to-go at turn $t$.
\end{itemize}
A held-out state never seen in training has no table entry; it gets the turn baseline (reported separately as
coverage).

\section{Setup}
Perfection, 16$\times$16 Continents, 3 players, difficulty Crazy (DECISION on \#166, 25 Nov), seat 0 Xin-Xi as
the game's human seat (\texttt{humans: [0]}, as on Jakob's branch \texttt{rl/learner-human}; on \texttt{dev}
the gate still has no human seat), opponents' tribes from the gate's rule, \texttt{greedy-full-probe} in every
seat. Training seeds $1, 2, \dots$; held-out seeds $500{,}000{,}000, \dots$ (10{,}000 games); both below the
gate's evaluation seeds ($\geq 10^9$). Training sizes 1k, 3k, 10k, 30k, 100k, 300k games (each a prefix of the
same seed list).

Commands:
\begin{verbatim}
cargo run --release -p rl-value -- collect --games 300000 --seed-start 1 \
  --out target/rl-value/train.tsv
cargo run --release -p rl-value -- collect --games 10000 --seed-start 500000000 \
  --out target/rl-value/test.tsv
cargo run --release -p rl-value -- fit --train target/rl-value/train.tsv \
  --test target/rl-value/test.tsv --sizes 1000,3000,10000,30000,100000,300000 \
  --out ../experiments/13-rl-full-state-value
\end{verbatim}
Sanity checks: unit tests (backward induction on two hand-made games; one row per learner turn from turn 0;
the same seed gives the same game). Pilot, 200 games: 0.9\,s on 10 threads, 30 learner turns per game, final
scores around 9{,}000.

\section{Results}
\ExpResults

\todo{claude}{Tables from value-summary.json: learning curve over training size; per-turn coverage and RMSE.}

\section{Behaviour}
\todo{claude}{Not a policy change: which states the bot visits (ranges, visits per state).}

\section{Interpretation}
\todo{claude}{Why the result came out this way, and whether the prediction held.}
\todo{david}{Your reading: agree, disagree, what is missing?}

\section{Threats to validity}
\todo{claude}{Two or three concrete threats.}

\section{Verdict and next step}
\todo{claude}{No gate run (no new player). What this means for value iteration.}

\section{Notes for the paper}
\todo{claude}{Draft: the one-sentence claim and which figure or table would carry it.}
\todo{david}{Is this worth a place in the talk or report? Where?}

\end{document}
